%-------------------------------------------------------------------------------
%	SECTION TITLE
%-------------------------------------------------------------------------------
\cvsection{Professional Experience}


%-------------------------------------------------------------------------------
%	CONTENT
%-------------------------------------------------------------------------------
\begin{cventries}
%---------------------------------------------------------
  \cventry
    {Senior NLP Engineer} % Job title
    {deepset} % Organization
    {Berlin, Germany} % Location
    {Feb 2024 - Present} % Date(s)
	{
      \begin{cvitems} % Description(s) of tasks/responsibilities
%
		% --- Maintainership -------------------------------------------------
		\item {Core maintainer of Haystack, an open-source LLM orchestration framework. Sustained day-to-day operations: new components and features, release cycles, pull request reviews, issue triage and mentoring external contributors. Shaped core abstractions: what components to add, their interfaces, how they interact with the pipeline core, and testing standards.}
%
		% --- Pipeline breakpoints & crash recovery --------------------------
		\item {Designed and shipped pipeline breakpoints, from proof-of-concept to core framework feature: pipelines and agents pause at any component, serialise state to a snapshot for inspection, and resume from it. Generalised the mechanism into crash recovery, so a failing run surfaces its last valid snapshot instead of a full restart.}
%
		% --- Retrieval & document store layer -------------------------------
		\item {Owned recurring work on the retrieval and document store layer: formalised the embedder and retriever interfaces as Protocols, added components against them (sentence-window and auto-merging retrievers, hybrid and multi-query retrieval, query expansion, recursive and embedding-based splitters), and consolidated the backends under one API and a shared test suite replacing per-backend duplicates. Benchmarked retrieval techniques and presented the results at PyCon Lithuania 2025.}
%
		% --- Telemetry & usage analytics ------------------------------------
		\item {Turned collected telemetry into the team's first reliable picture of adoption and component usage: set up an hourly export of raw events from PostHog into BigQuery, with SQL views, Looker Studio dashboards and automated weekly reports.}
%
		% --- n8n / Hayhooks integration -------------------------------------
		%\item {Prototyped and built the integration between Haystack and n8n, writing a custom TypeScript node supporting HTTP, MCP and A2A connections, with workflows indexing data from Gmail, Google Drive, and Notion, and an email-triggered question-answering assistant.}
%
        % --- Agents/Claude Code  -------------------------------------
		\item {Moved routine development tasks onto coding by writing and using agent skills for specific components, used during implementation, pull request reviews and release preparation. The effort shifted from writing code to specifying and verifying it.}
      \end{cvitems}
    }
	
%---------------------------------------------------------



  \cventry
    {Senior Data Scientist} % Job title
    {Veeva Systems} % Organization
    {Berlin, Germany} % Location
    {Mar 2023 - Jan 2024} % Date(s)
	{
      \begin{cvitems} % Description(s) of tasks/responsibilities
		\item {Built a PySpark MLlib classifier categorising life-science organisations from a 2 million-sample multilingual dataset combined with open data resources, cutting curator workload by 10\% within the Veeva Link product.}
		\item {Improved the researcher-to-activity linking model with pre-trained static embeddings over a 3.2 billion-sample annotated dataset, for a 2\% gain in model performance.}
		% \item {Implemented Python development best practices (tests, linting, PEP 8, type hints) within an existing CI pipeline and mentored junior data scientists in applying the same practices.}
		\item {Owned the Key Account newsfeed relevance model end to end: data collection, annotation batches, inter-annotator agreement, and deployment as an API for the Data Engineering team; used contrastive learning on pre-trained sentence-transformers to handle a corpus with only roughly 5\% relevant samples.}
      \end{cvitems}
    }
	
	% This experience taught me valuable lessons in data engineering, optimization, and scalability. I gained expertise in optimizing PySpark workflows, efficiently handling data partitioning, and utilizing distributed computing resources to ensure seamless model training and evaluation. Overall, my experience at Veeva allowed me to not only work with massive datasets but also to hone my skills in data engineering, machine learning.
	
	% Improved the Key Account newsfeed with a model developed identify relevant news, this was challenging due to imbalanced dataset only roughly 5\% of the news articles are relevant and not having access to the full news article text, I used a contrastive learning approach based on the sentence-transformers. I was responsible for collecting raw data, creating annotation batches, measuring the inter-annotation agreement and give feedback to the anntoation team, and deploying into production the model as an API consumed by the Data Engineering team.

%---------------------------------------------------------
  \cventry
    {Lead NLP Engineer} % Job title
    {TripActions (now Navan)} % Organization
    {Berlin, Germany} % Location
    {May 2022 - Feb 2023} % Date(s)
	{
      \begin{cvitems} % Description(s) of tasks/responsibilities
		\item {Joined through TripActions' acquisition of Comtravo; the team was tasked with building the NLP/NLU component of a customer-support chatbot, carrying over components previously developed for the email channel.}
		\item {Together with the team defined an intent and entities schema, set up the infrastructure for annotation based on the argilla.io annotation tool and onboarded annotators.}
		\item {Defined and implemented an entropy-based active-learning strategy to select samples for annotation.}
		\item {Fine-tuned, evaluated and threshold-tuned classification and sequence-tagging models on the collected annotated data, based on pre-trained Transformer architectures.}
      \end{cvitems}
    }


%---------------------------------------------------------

\cventry
    {Lead NLP Engineer} 
    {Comtravo}
    {Berlin, Germany}
    {May 2021 - Feb 2022}
    {
      \begin{cvitems} 
		  \item {Led a team of 3 developers guiding them in technical decisions and supervised system changes, and also managed the 4 annotators.}
		  \item {Managed the bi-weekly sprint planning and execution, coordinating development tasks based on system performance and feature requests.}
		  \item {Defined quantifiable measures in collaboration with the Data Engineering team, resulting in performance reports and monitoring dashboards.}
		  \item {Supervised the data annotation by ensuring annotation quality and consistency across the whole corpora.}
		  \item {Continued to maintain an active role in software development: implementing new features, maintenance, refactoring and code reviews.}
        \end{cvitems}
 	}


  \cventry
    {Senior NLP Engineer}
    {Comtravo} % Organization
    {Berlin, Germany} % Location
    {Aug 2017 - Apr 2021} % Date(s)
    {
      \begin{cvitems}
		\item {Helped build an automation system from its inception, turning incoming business-travel emails into trip itinerary offers for travel agents.}
		\item {Trained, evaluated and deployed supervised models for text classification and sequence tagging using different architectures.}
		\item {Designed and constructed Knowledge Bases containing worldwide descriptions of airport and train stations based on open resources and in-house operational data.}
		\item {Conceived and implemented Entity-Linking approaches to map entities (e.g: airports, train stations, hotels) recognised in incoming travel requests into unique identifiers in Knowledge Bases.}
		\item {Collaborated together with the team in developing a ranking and selection algorithm of trip itineraries offers based on the Pareto Efficiency.}
		\item {Developed several modularised Python components with type-annotations, linting and test coverage of 95\%.}
		\item {The automation system generating travel itineraries from customer emails reduced by 30\% the workload of travel-agents.}
        \end{cvitems}
 	}

%------------------------------------------------------------------------------------------------------------


\begin{comment}
  \cventry
    {Lead NLP Engineer} % Job title
    {Comtravo GmbH} % Organization
    {Berlin, Germany} % Location
    {Aug 2017 - Apr 2022} % Date(s)
	{
      \begin{cvitems} % Description(s) of tasks/responsibilities
		\item {Joined as an NLP Engineer in August 2017, was promoted to Senior NLP Engineer in July 2019, and in June 2021 to Lead NLP Engineer.}
        \item {Led a team of 3 developers + 4 annotators, working on the system that automatically answers incoming email travel requests and assists travel agents in handling them. Coordinating development tasks based on system performance and feature requests.}
        \item {Trained, evaluated and improved different models for text classification and fine-grained NER, increasing the performance of identifying specific booking requests and performing information extraction to automatically fulfil booking requests.}
		\item {Developed algorithms to map input text into unique Knowledge Base identifiers, e.g: airports, train stations, hotels, geographic locations.}		
      \end{cvitems}
    }
\end{comment}


%---------------------------------------------------------
  \cventry
    {Data Engineer} % Job title
    {HelloFresh} % Organization
    {Berlin, Germany} % Location
    {Jan 2016 - Jun 2017} % Date(s)
    {
      \begin{cvitems} % Description(s) of tasks/responsibilities
        \item {I was part of the initial Data Engineering team tasked with transitioning from crontab-based processes to an ETL platform.}
        \item {Developed the initial prototype for ETL management using Airflow, which evolved into the production data pipeline management system.}
		%\item {Conceptualised and developed multiple ETL processes with the PySpark ecosystem and orchestrated by Airflow.}
        \item {Developed a classifier to identify customer feedback review mentioning issues with the ordered meal kits.}
      \end{cvitems}
    }
	
%---------------------------------------------------------
  \cventry
    {Researcher and Developer} % Job title
    {INESC-ID, Instituto Superior Técnico} % Organization
    {Lisbon, Portugal} % Location
    {Jun 2011 - Apr 2014} % Date(s)
    {
      \begin{cvitems} % Description(s) of tasks/responsibilities
        \item {\href{http://arquivo.pt/wayback/20151118124735/http://dmir.inesc-id.pt/project/Reaction}{REACTION - Computational Journalism}: a research project on extracting and connecting entities, relations and topics from news articles.}
        \item {Explored and implemented methods for entity-relationship extraction, based on: hand-built patterns, supervised linear classifiers with linguistic features and semi-supervised methods based on seed relationships and large amounts of unannotated text.}
		\item {Developed an entity-linking approach associating personalities in news articles with Wikipedia using textual similarities and its graph structure.}
        \item {Built a graph based on topic models extracted from news articles and person co-occurrences, allowing to explore topics connecting persons.}
      \end{cvitems}
    }


%---------------------------------------------------------
  \cventry
    {Researcher and Developer} % Job title
    {LaSIGE, University of Lisbon} % Organization
    {Lisbon, Portugal} % Location
    {Sep 2008 - Oct 2010} % Date(s)
    {
      \begin{cvitems} % Description(s) of tasks/responsibilities
        \item {\href{https://arquivo.pt/wayback/20181014115718/http://xldb.di.fc.ul.pt/wiki/Grease}{GREASE - Geographic Reasoning for Search Engines}: a research project on resolving place names and reasoning over geographic ontologies.}
        \item {Developed a method to disambiguate toponyms based on textual context, information content and topological similarity measures.}
        \item {Built a method to align two geo-ontologies into a single linked-data ontology, exposing features from both to a single SPARQL query.}
		% \item {Published scientific articles presenting the results of these tasks.}
		\item {Ran n-gram language identification over a 10 million-page web crawl, then mined top Portuguese names overlapping with toponyms.}
      \end{cvitems}
    }

%---------------------------------------------------------
  \cventry
    {Software Developer} % Job title
    {Nokia Siemens Networks} % Organization % - \href{https://arquivo.pt/wayback/20181014115718/http://xldb.di.fc.ul.pt/wiki/Grease}{Geographic Reasoning for Search Engines}
    {Lisbon, Portugal} % Location
    {Oct 2007 - Jul 2008} % Date(s)
    {
      \begin{cvitems} % Description(s) of tasks/responsibilities
        \item {Developed data collection modules for a GSM monitoring system using Java, CORBA Architecture and Oracle RDBMS.}
		% \item {Held as an 80\% position, with the remaining 20\% dedicated to completing the M.Sc. curriculum component.}
      \end{cvitems}
    }

\end{cventries}
